% ===== BOT 11: Hàm mất mát của SADGS — L1, SSIM, tone-curve, frequency loss =====

\begin{frame}{Hàm mất mát của SADGS: $L_1$ + SSIM + $L_2$ (và hai loss chưa dùng)}
\scriptsize
\begin{columns}[T]
\begin{column}{0.47\linewidth}
\textbf{Loss thực sự được tối ưu} --- \texttt{train.py:259}:
\[
\mathcal{L} = (1-\lambda_{\text{dssim}})L_1 + \lambda_{\text{dssim}}(1-\mathrm{SSIM}) + \lambda_{L_2} L_2
\]
\[
= \;0.8\,L_1 \;+\; 0.2\,(1-\mathrm{SSIM}) \;+\; \mathbf{2.0}\,L_2
\]
$\lambda_{\text{dssim}}=0.2$, $\lambda_{L_2}=2.0$ (\texttt{arguments/\_\_init\_\_.py:86,117}). SSIM dùng kernel CUDA \texttt{fused\_ssim} (\texttt{train.py:18,258}), không dùng \texttt{ssim()} thuần torch ở \texttt{loss\_utils.py:46} (cửa sổ Gaussian $11\times11$, $\sigma=1.5$, $C_1=0.01^2$, $C_2=0.03^2$).
\vspace{3pt}

\alert{Hai loss được định nghĩa nhưng KHÔNG vào $\mathcal{L}$:}
\begin{itemize}
    \item \texttt{tone\_curve\_loss} (\texttt{:28}): $\;\dfrac{1}{N}\sum\Bigl(\dfrac{p-g}{\mathrm{sg}(p)+\varepsilon}\Bigr)^{2}$, $\varepsilon=10^{-6}$ --- chia sai số cho độ sáng dự đoán (stop-gradient) $\Rightarrow$ học trên \emph{sai số tương đối}, khuếch đại gradient vùng tối.
    \item \texttt{frequency\_loss} (\texttt{:80}): $\;\dfrac{1}{N}\sum \mathrm{ReLU}\bigl(\sigma_{\text{lin}} - \lambda_{\text{lim}}\bigr)$ với $\sigma_{\text{lin}}=\sqrt{\sqrt{\det\Sigma_{2D}}}$ và $\lambda_{\text{lim}} = 1/\sqrt{S_{xx}+S_{yy}}$ lấy mẫu từ \texttt{st\_map} tại \texttt{means2D} --- phạt Gaussian \textbf{to hơn bước sóng} texture.
\end{itemize}
Cả hai có trọng số $=0$ và \textbf{không được tham chiếu} ở bất kỳ đâu (\texttt{lambda\_tone}, \texttt{lambda\_freq} --- \texttt{args:118-119}).
\end{column}
\begin{column}{0.51\linewidth}
\centering
\includegraphics[width=\linewidth]{sadgsx/11_frequency_loss.png}
\captionof{figure}{\scriptsize \texttt{frequency\_loss} (\texttt{loss\_utils.py:80-110}): bước sóng cho phép $\lambda_{\text{lim}}$ co lại ở vùng texture rậm; Gaussian có $\sigma_{\text{lin}}$ vượt ngưỡng bị phạt tuyến tính một phía.}
\end{column}
\end{columns}
\vspace{2pt}
\begin{itemize}
    \item[$\Rightarrow$] Loss của SADGS $=$ loss 3DGS gốc \textbf{$+$ một số hạng $L_2$ trọng số lớn}. Cơ chế "structure-aware" nằm ở \textbf{densification} (thống kê $\eta$ không gradient, \texttt{freq\_utils.py:181}) chứ không nằm trong hàm mất mát.
\end{itemize}
\end{frame}
